% discussion.tex -- BNR-2026-013

\section{Discussion}
\label{sec:discussion}

\subsection{Proving Time Convergence}

The most significant finding is the scale-dependent convergence
of halo2-KZG proving time toward Groth16. At 10~arithmetic
steps, halo2-KZG is 4.7$\times$ slower; at 500~steps, the ratio
drops to 1.2$\times$. This convergence is explained by three
factors:

\begin{enumerate}
  \item \textbf{Fixed overhead amortization}: halo2-KZG incurs
    fixed costs for FFT setup, polynomial commitment, and SRS
    loading that are independent of circuit size. These costs
    dominate at small scales but become negligible as row count
    increases.
  \item \textbf{Equivalent asymptotic complexity}: Both Groth16
    and halo2-KZG have $O(n \log n)$ proving complexity, so
    per-row marginal cost is comparable.
  \item \textbf{Power-of-two padding}: halo2 circuits are padded
    to the next power of two ($2^k$ rows). Small circuits pay
    significant padding overhead (e.g., 11 constraints padded to
    32~rows), but this penalty diminishes at larger scales.
\end{enumerate}

At production scale---enterprise state transition circuits with
30K--50K rows after custom gate optimization---the proving time
ratio is projected to approach 1.0$\times$ or potentially favor
halo2-KZG due to reduced row count.

\subsection{Custom Gate Impact}

The hash chain benchmark directly measures the impact of custom
gates: the degree-5 S-box gate ($s \cdot (a^5 + c - b) = 0$)
reduces 3~R1CS constraints to 1~PLONKish row per hash round,
a 2.4$\times$ row reduction at chain length~100 (301~constraints
vs 128~rows).

For full Poseidon hashing (width~3, 8~full rounds + partial
rounds), the projected reduction is more dramatic: $\sim$211~R1CS
constraints reduce to $\sim$12 halo2 rows using degree-5 S-box
gates, representing a 17$\times$ reduction. Since the Basis
Network state trie is Poseidon-based, this reduction translates
directly to faster proving.

The current Validium circuit (274,291 R1CS constraints) is
projected to reduce to 30K--50K halo2 rows with full custom gate
optimization (ArithmeticGate, PoseidonGate, BitwiseGate,
StorageGate), representing a 5--9$\times$ circuit size reduction.

\subsection{Proposed Custom Gate Architecture}

Table~\ref{tab:gate-arch} presents the proposed custom gate
architecture for the Basis Network zkEVM prover, informed by
Scroll and zkSync patterns.

\begin{table}[htbp]
\caption{Proposed Custom Gate Architecture for zkEVM}
\label{tab:gate-arch}
\centering
\begin{tabular}{@{}lll@{}}
\toprule
\textbf{Gate} & \textbf{EVM Opcodes} & \textbf{Reduction} \\
\midrule
ArithmeticGate & ADD, SUB, MUL, DIV & 100--250$\times$ \\
BitwiseGate & AND, OR, XOR, SHL, SHR & 25--31$\times$ \\
ComparisonGate & LT, GT, SLT, SGT, EQ & 8--50$\times$ \\
MemoryGate & MLOAD, MSTORE, MCOPY & Permutation \\
StorageGate & SLOAD, SSTORE & Poseidon + path \\
StackGate & PUSH, POP, DUP, SWAP & Permutation \\
PoseidonGate & State root computation & 600--900$\times$ \\
HashGate & KECCAK256 & 5--10$\times$ \\
\bottomrule
\end{tabular}
\end{table}

\subsection{Migration Strategy}

We propose a two-phase migration to minimize risk:

\textbf{Phase~1: Dual Verification.} Deploy a PLONKVerifier
contract alongside the existing Groth16Verifier on the Basis
Network L1. The BasisRollup contract accepts both proof types
via a router pattern. New enterprise chains deploy with PLONK;
existing chains continue Groth16. This ensures zero downtime
and provides a fallback path.

\textbf{Phase~2: PLONK-Only.} After validation, all enterprise
chains migrate to the halo2-KZG prover. The Groth16Verifier
remains deployed for historical proof verification but is
deprecated for new batches. The universal SRS is published and
frozen for the deployment.

\subsection{SRS Requirements}

The universal Structured Reference String must support circuits
up to a maximum row count:

\begin{itemize}
  \item Current Groth16 circuit: 274K R1CS constraints
  \item Enterprise circuits: $\sim$500 constraints/tx $\times$
    1,000~tx/batch = $\sim$500K rows
  \item Safety margin: $2^{20} \approx 1$M rows ($k = 20$)
\end{itemize}

PSE's converted perpetual-powers-of-tau files provide SRS for
$k$ up to 26 ($\sim$67M rows). Scroll uses $k \in \{20, 24, 26\}$.
For Basis Network, $k = 20$ is sufficient for initial deployment,
with headroom to $k = 24$ for future circuit growth.

\subsection{Projected Production Performance}

Based on our benchmarks, published Axiom data, and custom gate
analysis, we project the following production metrics for the
Basis Network zkEVM L2 prover:

\begin{itemize}
  \item \textbf{Row count}: 30K--50K rows (5--9$\times$ reduction
    from 274K R1CS constraints via custom gates)
  \item \textbf{Circuit size}: $k = 16$--17 (64K--128K rows with
    power-of-two padding)
  \item \textbf{Proving time}: 2--8~s (based on Axiom's 2~s for
    ECDSA at $k = 15$)
  \item \textbf{Proof size}: 672--900~bytes
  \item \textbf{Verification gas}: 290--420K ($< 500$K budget)
  \item \textbf{Setup}: Universal SRS (one-time, reusable)
\end{itemize}

\subsection{FFLONK Alternative}

For applications requiring minimum verification gas, a two-layer
approach---halo2-KZG inner proof with FFLONK outer wrapper---could
reduce L1 verification to $\sim$200K gas (the Polygon zkEVM
pattern). However, on the zero-fee Basis Network L1, the
additional 120--220K gas savings do not justify the implementation
complexity. The simpler single-layer halo2-KZG at 290--420K gas
is recommended.

\subsection{Risk Assessment}

Table~\ref{tab:risks} identifies migration risks and mitigations.

\begin{table}[htbp]
\caption{Risk Assessment}
\label{tab:risks}
\centering
\begin{tabular}{@{}p{2.5cm}lp{3cm}@{}}
\toprule
\textbf{Risk} & \textbf{Severity} & \textbf{Mitigation} \\
\midrule
Verification gas exceeds 500K & Low & Axiom production: 420K \\
SRS ceremony compromise & Medium & PSE PoT (71+ participants) \\
Axiom fork unmaintained & Medium & PSE fork fallback; stable API \\
Custom gate design errors & High & Coq verification; MockProver \\
Proving time regression & Low & Poseidon-heavy = PLONK advantage \\
\bottomrule
\end{tabular}
\end{table}

\subsection{Limitations}

Our benchmarks use small circuits (10--500 steps) on a commodity
desktop. While the convergence trend is clear, production-scale
validation at 30K--50K rows on server hardware is needed before
deployment. Published production data from Scroll, Axiom, and
Taiko corroborates our projections but uses different circuit
structures. The hash chain benchmark approximates Poseidon via
an $x^5$ S-box; a full Poseidon circuit would provide more
precise constraint reduction measurements.
